% GENERADO AUTOMATICAMENTE por research/eval/report/tabla_arquitecturas.py
% No editar a mano.
\begin{table}[h]
\centering
\footnotesize
\setlength{\tabcolsep}{4pt}
\resizebox{\textwidth}{!}{%
\begin{tabular}{|l|l|r|r|r|r|r|r|c|l|}
    \hline
    \textbf{Corpus} & \textbf{Modelo} & \textbf{Par. (M)} & \textbf{WER} &
    \textbf{Crít.} & \textbf{Crít. s/n} & \textbf{Vel.} &
    \textbf{$\Delta$ (pp)} & \textbf{IC 95\%} & \textbf{Veredicto} \\
    \hline
    \texttt{teleconciencia-es-400} & \texttt{whisper-medium} & 764 & 19.55 & 22.17 & 23.08 & 8$\times$ & --- & --- & referencia \\
    \texttt{teleconciencia-es-400} & \texttt{parakeet-tdt-0.6b-v3} & 627 & 19.83 & 10.36 & 10.00 & 106$\times$ & +0.28 & [-1.28, +2.03] & sin evidencia \\
    \texttt{voxpopuli-es-400} & \texttt{whisper-medium} & 764 & 9.63 & 17.52 & 7.74 & 8$\times$ & --- & --- & referencia \\
    \texttt{voxpopuli-es-400} & \texttt{parakeet-tdt-0.6b-v3} & 627 & 6.83 & 4.46 & 4.33 & 44$\times$ & -2.80 & [-3.48, -2.20] & mejora \\
    \hline
\end{tabular}}
\caption{Modelo base de la comparativa frente a una arquitectura de transductor, sobre
los mismos clips. \textbf{Crít.} es la tasa de error sobre negaciones, numerales y
cuantificadores, y \textbf{Crít. s/n} la misma tasa excluyendo numerales, cuyo recuento
mide en parte la convención de escritura: la referencia escribe «dos mil diez» donde
Whisper escribe «2010», y el alineamiento contabiliza tres errores por una cifra bien
reconocida. \textbf{Vel.} es cuántas veces más rápido que el audio se procesa.
$\Delta$ negativo indica mejora respecto al modelo base. Las decodificaciones no son
equiparables entre arquitecturas: cada sistema se evalúa en su configuración estándar.}
\label{tab:comparativa-arquitecturas}
\end{table}
